\documentclass[11pt,a4paper]{article}
\usepackage[T1]{fontenc}
\usepackage[utf8]{inputenc}
\usepackage{lmodern}
\usepackage[margin=27mm]{geometry}
\usepackage{amsmath,amssymb,amsthm}
\usepackage{booktabs,array}
\usepackage{microtype}
\usepackage{xurl}
\usepackage[hidelinks]{hyperref}
\usepackage{fancyvrb}
\hypersetup{pdftitle={A conditional refinement of the Subset Sum exponent to 0.487676},pdfsubject={Parameter refinement and transfer of checksum certificates}}
\newtheorem{theorem}{Theorem}
\newtheorem{lemma}[theorem]{Lemma}
\newtheorem{proposition}[theorem]{Proposition}
\theoremstyle{remark}
\newtheorem*{remark}{Remark}
\newcommand{\Base}{\ensuremath{\mathcal{B}}}
\newcommand{\poly}{\operatorname{poly}}
\newcommand{\Torus}{\mathbb{T}}
\newcommand{\eps}{\varepsilon}
\newcommand{\logtwo}{\log_2}
\setlength{\emergencystretch}{2em}
\renewcommand{\arraystretch}{1.16}
\allowdisplaybreaks
\title{A conditional refinement of the Subset Sum\texorpdfstring{\\}{ }exponent to \texorpdfstring{$0.487676$}{0.487676}}
\author{}
\date{20 October 2026}

\begin{document}
\maketitle

\begin{abstract}
We give a parameter and certificate refinement of the claimed
$O(2^{0.4877n})$ Subset Sum construction in the pinned manuscript
\Base~\cite{baseline}. The internal work exponent is reduced from
$0.48769$ to $0.48767595$. A transfer argument preserves the existing
43 filters and derives the required new moment bounds from their
430 original certificates, without new quadrature. Together with exact
rational checks of the alias constraints, this gives the conditional
bound $O_c(2^{0.487676n})$ for every fixed polynomial input-bit-length
bound $b\le n^c$. The qualification is essential: the algorithmic lemmas
and original moment inequalities of \Base\ are inherited premises, not
independently verified here. This note is not a formal proof-assistant
verification, a literature-priority claim, or a practical speedup claim.
\end{abstract}

\section{Statement and dependencies}
\label{sec:statement}

Let \Base\ denote \texttt{subset-sum/PROOF\_0.4877.tex} in
\texttt{a-kapila/maths-proofs}, pinned at commit
\texttt{44c34961b56784c220a6e913cc068d57dc37c5a3}~\cite{baseline}.
That manuscript relies on the 4 October 2026 preprint
\emph{Subset Sum in Time $O(2^{0.49n})$}~\cite{original}.
The intended specification is the Subset Sum benchmark~\cite{benchmark};
no assertion about its current leaderboard is needed for this note.

The inherited premises are the isolation and implementation lemmas,
the preliminary-search and alias arguments under their numerical side
conditions, the parameterised checksum construction, finite-circle
convergence, and the certified continuous moment bounds of \Base.
References to numbered statements of \Base\ refer to that pinned version.

\begin{theorem}[Conditional refinement]
\label{thm:main}
Assume these inherited premises of \Base\ are valid. There is a single
uniform randomised classical word-RAM program which halts on every input
and is correct with probability at least $2/3$. For each fixed positive
integer $c$, there are constants $C_c,N_c>0$ such that, for every input
with $n\ge N_c$ and $b\le n^c$, every random execution uses at most
\[
 C_c\,2^{0.487676n}
\]
word operations. Here $b$ is the largest bit length of an input integer
or the target, with the conventions of \Base, and the word length is
\[
 w_{\mathrm{RAM}}=\left\lceil4\bigl(n+b+\logtwo(n+2)\bigr)\right\rceil.
\]
The program does not depend on $c$.
\end{theorem}

The baseline internal constants are
\begin{equation}
\label{eq:baseline}
 \tau_b=0.48769,\qquad \eta_b=0.51613,\qquad
 \Delta_b=0.02844,\qquad \Gamma_b=0.45925.
\end{equation}
Its displayed exponent $0.4877$ absorbs overhead beyond the linear
exponent $0.4876902$. The refinement below changes the internal work
exponent; it does not merely spend the baseline's rounding slack.

\section{Parameters and input guard}
\label{sec:parameters}

All terminating decimals in the parameter table are exact rational
numbers. The machine-readable values are in \texttt{parameters.json}.
\begin{center}
\begin{tabular}{@{}lll@{}}
\toprule
Parameter & Exact value & Meaning\\
\midrule
$\tau$ & $0.48767595$ & Internal work exponent\\
$\eta$ & $0.51612858$ & Large-prime exponent\\
$a$ & $0.008145205148$ & Mask fraction\\
$\mu$ & $0.475442899283$ & Large-block fraction\\
$\gamma$ & $0.060035950245$ & Small-prime parameter\\
$l$ & $0.483780402235$ & Left-block fraction\\
$r=1-\mu-l$ & $0.040776698482$ & Right-block fraction\\
$\eps$ & $10^{-11}$ & Pattern accuracy\\
$\zeta$ & $10^{-11}$ & Preliminary-search slack\\
\bottomrule
\end{tabular}
\end{center}
The new checksum budgets are
\begin{equation}
\label{eq:budgets}
 \Delta=\eta-\tau=0.02845263,\qquad
 \Gamma=2\tau-\eta=0.45922332.
\end{equation}
Retain the baseline's 43 filters, filter-selection rule, isolation
transformation, and final decision rule. Construct checksum tables and
choose the sampling fraction using the \emph{new} budgets. The regions
used in the certificate analysis will retain the \emph{old} boundaries.

Replace the input guard by
\begin{equation}
\label{eq:guard}
 10^{11}\left\lceil\logtwo(b+1)\right\rceil\le n.
\end{equation}
When the guard fails, or $n$ is below a sufficiently large fixed cutoff,
use the exact meet-in-the-middle fallback of \Base. For each fixed $c$,
the guard eventually holds whenever $b\le n^c$. On inputs passing it,
\begin{equation}
\label{eq:bitfactor}
 b+2\le 2^{10^{-11}n+1}.
\end{equation}
The bit-length factor in the alias analysis therefore consumes at most
$10^{-11}$ of the linear-exponent budget.

\section{Alias work and coverage}
\label{sec:alias}

Write
\[
 h(x)=-x\logtwo x-(1-x)\logtwo(1-x).
\]
The strict alias constraints of \Base, Section 2.2, become
\begin{align}
 \gamma &< \tfrac12 h(2a), \tag{C1}\label{c1}\\
 l+\bigl(h(a)-\gamma\bigr)\mu &<\tau, \tag{C2}\label{c2}\\
 r+(1-\gamma)\mu &<\tau, \tag{C3}\label{c3}\\
 1-\mu+\bigl(h(a)+h(\tfrac12-a)-\gamma\bigr)\mu-\eta
 &<\tau, \tag{C4}\label{c4}\\
 h(\tfrac12-a)\mu+\tfrac12-\tau &<\tau, \tag{C5}\label{c5}\\
 0<r<\tau,\qquad l&<\tau, \tag{C6}\label{c6}\\
 \mu(1-2a)&<l, \tag{C7}\label{c7}\\
 \tfrac{\mu}{2}\logtwo3&<\eta. \tag{C8}\label{c8}
\end{align}
These respectively control representation coverage, left lists, right
value classes, matching, mask enumeration, value-map construction,
deficiency accounting, and reduction on half-blocks.

The accompanying exact-rational interval checker verifies all eight
constraints. The tightest five have the following approximate slacks:
\begin{center}
\begin{tabular}{@{}lr@{}}
\toprule
Constraint & Positive slack\\
\midrule
Coverage, (C1) & $1.84806989556\times10^{-8}$\\
Left list, (C2) & $1.84818109678\times10^{-8}$\\
Right value classes, (C3) & $1.84806927342\times10^{-8}$\\
Matched pairs, (C4) & $1.84813377432\times10^{-8}$\\
Right-mask enumeration, (C5) & $1.84815267754\times10^{-8}$\\
\bottomrule
\end{tabular}
\end{center}
Every tested alias and auxiliary size inequality has slack at least
\begin{equation}
\label{eq:alias-slack}
 \delta_{\mathrm{alias}}=1.8480692734174335\times10^{-8}.
\end{equation}
The auxiliary checks include
\begin{equation}
\label{eq:aux}
 0<a<\tfrac14,\qquad
 \mu(1-2a)<\tfrac12,\qquad
 \tfrac12h(2a)+h(a)<1-2a,
\end{equation}
and $h(a)<h(\tfrac12-a)$. The layer inequality supplies the combinatorial
size condition in the averaged small-layer argument of \Base.

In particular,
\[
 \zeta+10^{-11}=2\times10^{-11}<\delta_{\mathrm{alias}}.
\]
Thus the preliminary-search slack and input-bit-length factor fit
strictly inside the alias work margin. Under the inherited alias
premises, the preliminary-search, coverage, expected-work and cap
arguments carry over with the new constants. The capped alias cost is
\begin{equation}
\label{eq:alias-time}
 2^{\tau n+O(n^{0.7})}\poly(n+b).
\end{equation}
Strict inequalities absorb the integer rounding of block and mask
sizes for sufficiently large $n$, as in \Base.

\section{Transferring the checksum certificates}
\label{sec:transfer}

Optimising the alias constraints alone would not establish the claimed
refinement. The checksum budgets change as well, so both moment bounds
and their domination of sampled patterns must be preserved.

Use the baseline score functions $u_v(x),w_v(x)$ for
$v\in\{0,\ldots,42\}$. For each pattern, the selected filter maximises
its approximate score $U_v$; write $U,W$ for that filter's scores.
The baseline guarantees $W\ge0$. The new sampling fraction is
\begin{equation}
\label{eq:alpha}
 \alpha=\begin{cases}
 1,&W\le\Gamma,\\
 \Gamma/W,&W>\Gamma.
 \end{cases}
\end{equation}
A sampled pattern satisfies $\alpha U<\Delta$.

Set
\begin{equation}
\label{eq:transfer-params}
 q=\frac{\Gamma}{\Gamma_b},\qquad
 A_b=\frac{\Delta_b}{\Gamma_b},\qquad
 A=\frac{\Delta}{\Gamma},\qquad
 K=\frac{\Delta}{q}-\Delta_b.
\end{equation}
These constants satisfy
\begin{equation}
\label{eq:transfer-identity}
 0<q<1,\qquad K>0,\qquad A-A_b=K/\Gamma_b,
\end{equation}
with approximate values
\begin{align*}
 q&=0.999941905280348394120849\ldots,\\
 K&=0.00001428304359630517021652994\ldots,\\
 A-A_b&=0.00003110080260491055028095796\ldots.
\end{align*}
Retain the old region boundaries $\Gamma_b$ times
\[
 1,\ 1.25,\ 1.6,\ 2,\ 2.6,\ 3.5,\ 5,\ 8,\ 15,\ \infty.
\]
These are boundaries for the proof, not for constructing the new tables.

Each old certificate has a nonnegative multiplier vector $z$ and
an affine row vector $G_b(x)$. Write
\[
 I_b=\int_{\Torus} M(x)^2\,2^{z\cdot G_b(x)}\,dx,
 \qquad M(x)^2=2+2\cos x,
\]
where $\Torus=\mathbb{R}/2\pi\mathbb{Z}$ has normalised measure. Comparison
rows $u_v-u_j$ correspond to other filters $j$ and are nonnegative at
the pattern level because $v$ was selected to maximise $U_v$.

\begin{lemma}[Certificate transfer]
\label{lem:transfer}
Assume \eqref{eq:transfer-identity}. The old certificates can be modified
as follows. Each modification retains the required domination for
sampled patterns in the old region.
\begin{enumerate}
\item For $W\le\Gamma_b$, replace the first row $\Delta_b-u$ by
      $\Delta/q-u$, leaving all multipliers unchanged. Then
      $I'=2^{z_0K}I_b$.
\item For $\Gamma_bL<W\le\Gamma_bH$ with finite $H$ and $L\ge1$,
      replace $A_bw-u$ by $Aw-u$ and replace the multiplier of
      $\Gamma_bH-w$ by $z_2+z_0(A-A_b)$. Then
      $I'=2^{z_0KH}I_b$.
\item For $W>15\Gamma_b$, replace $A_bw-u$ by $Aw-u$ and leave all
      multipliers unchanged. If $w_v(x)\le B$ on the real circle, then
      $I'\le2^{Bz_0(A-A_b)}I_b$.
\end{enumerate}
Here $I'$ denotes the integral of the modified certificate.
\end{lemma}

\begin{proof}
In the low region, \eqref{eq:alpha} gives $q\le\alpha\le1$. Put
$\Theta=\Delta/q=\Delta_b+K$. Sampling implies
$U<\Delta/\alpha\le\Theta$. The old condition $z_0\ge1$ gives
\begin{equation}
\label{eq:low-domination}
 -\alpha U
 =-\Delta+\alpha(\Delta/\alpha-U)
 \le-\Delta+z_0(\Theta-U).
\end{equation}
All remaining region and comparison rows are nonnegative at the
pattern level. Pointwise, the modified exponent is the old exponent
plus $z_0K$, proving the first integral identity. This includes the
interval $\Gamma<W\le\Gamma_b$, where $\alpha$ is less than one.

In a bounded region the old first three rows are
\[
 \bigl(A_bw-u,\ w-\Gamma_bL,\ \Gamma_bH-w\bigr).
\]
The multiplier increase is nonnegative. The change in the total
exponent is exactly
\begin{equation}
\label{eq:band-cancellation}
 z_0(A-A_b)w+z_0(A-A_b)(\Gamma_bH-w)=z_0KH.
\end{equation}
The old condition $z_0L\ge1$ implies
$\alpha<q/L<1/L\le z_0$. Moreover $\alpha AW=\Delta$ and sampling
implies $AW-U>0$. Hence
\begin{equation}
\label{eq:band-domination}
 -\alpha U=-\Delta+\alpha(AW-U)
 \le-\Delta+z_0(AW-U).
\end{equation}
The remaining rows are again nonnegative. This proves the second case.

In the unbounded region, $15z_0\ge1$ and $\alpha<q/15$ give the same
domination argument. The exponent changes by $z_0(A-A_b)w(x)$, which
is at most $Bz_0(A-A_b)$. Positivity of $M(x)^2$ on the real circle
then proves the final integral bound.
\end{proof}

The lemma establishes ideal pattern-level domination. The effect of
approximating the coordinate scores is accounted for explicitly in
Section~\ref{sec:error}; it is not silently omitted from the transfer.

\subsection{A pointwise bound for the tail}

\begin{lemma}[Coarse score bound]
\label{lem:score}
For each of the 43 baseline filters, $|w_v(x)|\le256$ on the real circle.
\end{lemma}
\begin{proof}
The finite data in Appendix A of \Base\ give eight cosine coefficients
for each of $r_v,s_v$, each of absolute value at most $2$, and a power
parameter $1\le\kappa_v\le2$. Suppress the filter index. Thus
$|r(x)|,|s(x)|\le16$. With normalised integrals, put
\begin{align*}
 F&=2^r,& G&=2^s,& Z_f&=\int F,& Z_g&=\int G,\\
 Z_{f,\kappa}&=\int F^\kappa,&
 Z_{g,\kappa}&=\int G^\kappa.
\end{align*}
Let
\begin{align*}
 \mathcal P(x)&=\int F(x+y)G(y)\,dy,\\
 \mathcal Y(x)&=\int F(x+y)G(y)\bigl(r(x+y)+s(y)\bigr)\,dy.
\end{align*}
The identities inherited from \Base\ are
\begin{align*}
 L(x)&=\logtwo\mathcal P(x)-\logtwo(Z_fZ_g),\\
 B_*(x)&=\mathcal Y(x)/\mathcal P(x)-\logtwo(Z_fZ_g),\\
 H_*&=\logtwo Z_{f,\kappa}+\logtwo Z_{g,\kappa}
       -\kappa\logtwo(Z_fZ_g),\\
 w(x)&=(2-\kappa)B_*(x)-2L(x)+H_*.
\end{align*}
Positivity gives $|\logtwo(Z_fZ_g)|\le32$ and
$|\logtwo\mathcal P(x)|\le32$. Also $\mathcal Y/\mathcal P$ is a
weighted average of $r(x+y)+s(y)$, so its absolute value is at most
$32$. Consequently
\[
 |L|\le64,\qquad |B_*|\le64,\qquad |H_*|\le64\kappa.
\]
Because $2-\kappa\ge0$,
\[
 |w|\le(2-\kappa)64+128+64\kappa=256.
\]
\end{proof}

\section{Moment margins and approximation errors}
\label{sec:error}

The following rounded-down moment margins are inherited from
Proposition B and Section 8.2 of \Base. They are assumptions of this
note, not quadratures evaluated by the new checker.
\begin{center}
\begin{tabular}{@{}lc@{}}
\toprule
Old region & Lower bound on $\eta_b-\logtwo I_b$\\
\midrule
$\rho_0$: $W\le\Gamma_b$ & $0.00025$\\
$\rho_1$: $\Gamma_b<W\le1.25\Gamma_b$ & $0.00025$\\
$\rho_2,\ldots,\rho_8$ & $4$\\
$\rho_9$: $W>15\Gamma_b$ & $17$\\
\bottomrule
\end{tabular}
\end{center}
The finite source data give
\[
 \max_v z_0^{(v,\rho_0)}=17.211,\qquad
 \max_v z_0^{(v,\rho_1)}=8.4537.
\]
The baseline bounds every multiplier-vector sum by $Z=418.36$.
These values and the transcribed filter data are recorded in
\texttt{source\_facts.json}. The optional source audit checks the
transcriptions against a supplied local copy of the pinned TeX source.

\begin{proposition}[Transferred continuous moments]
\label{prop:moments}
Assuming the old moment margins above, every modified certificate has
$\eta-\logtwo I'>10^{-6}$. Every modified multiplier sum is less than
$419$.
\end{proposition}
\begin{proof}
Since $\eta-\eta_b=-0.00000142$, Lemmas~\ref{lem:transfer}
and~\ref{lem:score} give lower bounds
\begin{align*}
 m_0&=0.00025+(\eta-\eta_b)-17.211K
       >0.0000027545366639917,\\
 m_1&=0.00025+(\eta-\eta_b)-8.4537(1.25)K
       >0.0000976492929373937,\\
 m_{2:8}&=4+(\eta-\eta_b)-418.36(15)K
       >3.9103667682,\\
 m_9&=17+(\eta-\eta_b)-418.36(256)(A-A_b)
       >13.6690976448.
\end{align*}
All are strictly greater than $10^{-6}$. For the bounded regions use
$z_0\le Z$ and $H\le15$; for the tail use $z_0\le Z$.
The only multiplier increase is $z_0(A-A_b)$, so every modified sum
is at most $Z(1+A-A_b)<419$. These are exact rational inequalities
verified by \texttt{verify.py}.
\end{proof}

There are only finitely many new, fixed affine certificates. The
finite-circle convergence premise of \Base\ therefore supplies a
finite-circle moment margin of $10^{-6}$ for all sufficiently large
primes. The strict continuous margin in
Proposition~\ref{prop:moments} is what permits this invocation of
convergence; convergence itself is not proved by the scalar checker.

The coordinate-average errors inherited from the baseline are
$|U_v-\bar u_v|\le3\eps$ and $|W_v-\bar w_v|\le4\eps$.
Each modified affine row has error at most $6\eps$. In particular,
$Aw-u$ has error at most $(4A+3)\eps<4\eps$ since $A<1/4$;
comparison rows have error at most $6\eps$. Fixed region constants
introduce no additional error.

The domination proof and baseline mean-square-error argument therefore
leave the positive exponent
\begin{equation}
\label{eq:error-exponent}
 10^{-6}-(18+6\cdot419)\eps
 =9.7468\times10^{-7}.
\end{equation}
For each sampled pattern the corresponding mean-square error, restricted
to the baseline good-table event, is at most
\[
 2^{-(9.7468\times10^{-7})n+O(1)}.
\]
The $2^{O(\sqrt n)}$ patterns and polynomially many targets are absorbed
by this strict exponential decay. The inherited good-table probability,
Cauchy--Schwarz, Markov's inequality and a union bound give the same
constant absolute-error guarantees as \Base, with probability $1-o(1)$.

\section{Correctness and pure exponential time}
\label{sec:time}

\begin{proof}[Proof of Theorem~\ref{thm:main}]
Use the baseline program with the constants and guard of
Section~\ref{sec:parameters}, retaining its bounded samplers, caps and
exact fallback. The preliminary search and alias procedure remain valid
under the inherited premises and the strict checks in
Section~\ref{sec:alias}. Sections~\ref{sec:transfer}--\ref{sec:error}
supply the replacement checksum moment and error analysis.

The exact modular checksum minus the exact alias phase sum is the phase
sum over subsets representing the unreduced target. On a NO instance
it is zero. When isolation produces a unique representing subset for
some target, its absolute value is one. The threshold $1/2$ therefore
separates these cases on the inherited success events. Isolation and
the vanishing failure probabilities give correctness at least $2/3$
above a sufficiently large fixed cutoff. The fallback is exact below
the cutoff and whenever the guard fails.

The work caps and bounded samplers ensure halting and the resource bound
on every random outcome, not merely in expectation. Prime generation
costs $2^{\eta n/2}\poly(n+b)$, which is below the work budget. The total
cost has the inherited form
\begin{equation}
\label{eq:total-time}
 T(n,b)\le
 2^{(\tau+20\eps)n+O(n^{0.7})}\poly(n+b)
 =2^{0.4876759502n+O(n^{0.7})}\poly(n+b).
\end{equation}
For fixed $c$ and $b\le n^c$, the guard eventually holds and
\[
 0.487676-(\tau+20\eps)=4.98\times10^{-8}>0.
\]
This linear-exponent gap absorbs both the subexponential factor and
$\poly(n+n^c)$. Thus \eqref{eq:total-time} is
$O_c(2^{0.487676n})$. The inherited address-space argument applies with
the same word length. All cutoffs and program parameters are fixed
independently of $c$.
\end{proof}

\begin{remark}
The cutoff and hidden constants may be very large. This argument gives
an asymptotic conditional refinement, not a demonstrated practical
speedup. It does not establish that $0.487676$ is optimal even within
the inherited framework.
\end{remark}

\section{Reproducibility and verification scope}
\label{sec:verification}

The accompanying files are \texttt{paper.tex}, \texttt{verify.py},
\texttt{parameters.json}, \texttt{source\_facts.json}, and
\texttt{verification.txt}. A \texttt{Makefile} provides the convenience
targets \texttt{verify} and \texttt{pdf}; neither target accesses the
network. To repeat the scalar checks, run
\begin{Verbatim}[fontsize=\small]
python3 verify.py
\end{Verbatim}
An optional audit of the finite baseline data is available as
\begin{Verbatim}[fontsize=\small]
python3 verify.py --baseline /path/to/PROOF_0.4877.tex
\end{Verbatim}
The source file should be obtained from the exact commit in
reference~\cite{baseline}. This optional audit checks transcribed data
and multiplier rows; it does not check the original integrals or
algorithmic correctness, nor authenticate the downloaded file's
provenance.

The checker uses only Python's standard library. Every inequality
decision uses \texttt{fractions.Fraction}. Logarithms are enclosed by
range reduction and, for $1\le y\le2$, the expansion
\begin{equation}
\label{eq:log-series}
 \ln y=2\sum_{j=0}^{m-1}\frac{z^{2j+1}}{2j+1}+R_m,
 \qquad z=\frac{y-1}{y+1},
\end{equation}
with the exact rational enclosure
\begin{equation}
\label{eq:log-tail}
 0\le R_m\le\frac{2z^{2m+1}}{(2m+1)(1-z^2)}.
\end{equation}
The implementation uses $m=40$. Division by an interval for $\ln2$
gives intervals for $\logtwo$, and interval operations enclose the
entropy values. Displayed decimal values are rounded outwards and are
not used to decide whether inequalities hold.

The included run was repeated on 20 October 2026 and passed the new
alias constraints, auxiliary size conditions, guard margin,
moment-transfer bounds, multiplier-sum budget, error exponent and
pure-exponent gap. It uses the supplied finite-data transcriptions.
The optional full-source audit and the original 430 quadratures were
\emph{not} run as part of this verification. All original algorithmic
and analytic premises remain explicit dependencies.

\paragraph{Provenance and status.}
This research note and its accompanying checker were prepared with
ChatGPT. The new certificate-transfer argument is a written derivation;
the script checks its scalar side conditions. Neither the script nor
this note is a Lean verification or an independent verification of the
underlying Subset Sum algorithm. No global optimality or literature
priority is claimed.

\begin{thebibliography}{9}
\raggedright
\bibitem{baseline}
\texttt{a-kapila/maths-proofs},
\emph{Subset Sum in time $O(2^{0.4877n})$: a complete proof},
12 October 2026.
File \texttt{subset-sum/PROOF\_0.4877.tex}, pinned commit
\texttt{44c34961b56784c220a6e913cc068d57dc37c5a3}.
\url{https://github.com/a-kapila/maths-proofs/blob/44c34961b56784c220a6e913cc068d57dc37c5a3/subset-sum/PROOF_0.4877.tex}.
The title is the source's title; this note treats its results as
inherited premises.

\bibitem{original}
OpenAI,
\emph{Subset Sum in Time $O(2^{0.49n})$},
OpenAI Math Release preprint, 4 October 2026.
Repository entry pinned at commit
\texttt{fd4aeeb2ee4fc729c18d98444fed42fd0529eeeb}.
\url{https://github.com/openai/math/blob/fd4aeeb2ee4fc729c18d98444fed42fd0529eeeb/preprints/Subset-Sum-in-Time-2-power-0-49n-October-4-2026/README.md}.

\bibitem{benchmark}
Big O Speedrun, \emph{Subset Sum benchmark}.
\url{https://bigospeedrun.com/subset-sum}.
\end{thebibliography}
\end{document}
